\section{Hochschulsport}
Um abzuschalten, könnt ihr das Sportangebot der Uni nutzen. Jedes Semester -- und auch in der vorlesungsfreien Zeit -- werden zahlreiche Kampfsport-, Fitness-, und Tanzkurse sowie Mannschaftssportarten und sogar Mehrtagestouren angeboten. Ihr könnt auch exotische Sportarten wie etwa Quidditch oder Fechten ausprobieren. In einigen Bereichen gibt es sogar die Unterscheidung zwischen Anfängern und Fortgeschrittenen wie z.B. beim Tennis, somit kann wirklich jede teilnehmen, die sich für den Sport begeistert.

Der Großteil der Sportkurse finden im Neuenheimer Feld statt, mitzunehmen sind Studierendenausweis und meist saubere Hallenschuhe.

Die meisten Kurse sind kostenpflichtig, wenn auch sehr günstig (etwa \EUR{35} pro Kurs und Semester), ihr müsst euch zur Teilnahme über die Website des Hochschulsports\footnote{\url{https://hochschulsport.issw-hd.de/}} anmelden. Die Anmeldung findet jedes Semester am Wochenende vor Vorlesungsbeginn statt. Damit die Server nicht überlastet werden, sind die Anmeldezeiten gestaffelt. Die genaue Uhrzeit, zu der ihr euch zu euren gewünschten Kursen anmelden könnt, findet ihr online bei der jeweiligen Kursbeschreibung. Wenn ihr einen Platz wollt solltet ihr, wenn die Anmeldung eröffnet ist, von der Liste der Kurse aus beim gewünschten Kurs auf Info und Buchung klicken, dann ist unten auf der Seite der Buchungsknopf. Bei manchen Kursen sind alle Plätze innerhalb der ersten Minuten vergeben. Um dann tatsächlich in den Kurs zu können müsst ihr noch am Tag an dem dieser stattfindet ein Teilnahmeticket erstellen, dann erhaltet ihr einen QR-Code mit dem ihr dann in das Gebäude gelassen werdet.

